%%%%%%%%%%%%%%%%%%%%%%%%%%%%%%%%%%%%%%%%%%%%%%%
%%      Pretty preamble for Astrophysics     %%
%%  (c) Jakob Lysgaard Rørsted (Mosumgaard)  %%
%%                                           %%
%%              Made for luaLaTeX            %%
%%%%%%%%%%%%%%%%%%%%%%%%%%%%%%%%%%%%%%%%%%%%%%%

%%%%%%%%%%%%%%%%%%%%%%%%%%%%%%%%%%%%%%%%%%%%%%%
% Core
%%%%%%%%%%%%%%%%%%%%%%%%%%%%%%%%%%%%%%%%%%%%%%%
% Essentials
\usepackage[final]{microtype}

% Language
\usepackage{babel}


%%%%%%%%%%%%%%%%%%%%%%%%%%%%%%%%%%%%%%%%%%%%%%%
% Fonts -- definiton from thesis (overkill?)
%%%%%%%%%%%%%%%%%%%%%%%%%%%%%%%%%%%%%%%%%%%%%%%
%  *** Comments regarding math font selection ***
%
%  Use the main font as math font for the glyphs it has. Use a fallback for
%  the rest. \setmathfont with range={…} and (s)script-features={…} are from
%  unicode-math. Script=Latin is directly from fontspec.
%
%  Script=Latin tells fontspec to use the OpenType script 'Latin' instead
%  of the expected 'Math'. The 'Math' script is not defined in normal OpenType
%  fonts as they are not math types.
%
%  (s)script-features={} refer to script unrelated the OpenType script above.
%  It overwrites unicode-maths normal OpenType feature flags for super- and
%  subscript. It defaults to '+ssty', which is also missing from the fonts.

% Use fontspec to set the main (OpenType) fonts
\usepackage{fontspec}

% Use the unicode math package to enable the use of OpenType math fonts
% --> Silence warnings related to mathtools
\usepackage[warnings-off={mathtools-colon,mathtools-overbracket}]{unicode-math}
% \usepackage{unicode-math}
\unimathsetup{math-style=ISO,bold-style=ISO,partial=upright}

% Use \lining to get lining numbers in a scope.
\newcommand{\lining}{\addfontfeature{Numbers=Lining}}

% Supress the 'Th' ligature (by default active in Arno Pros 'liga' feature). We
% create a new OpenType feature called 'noth' that substitute the 'T_h' ligature
% with 'T' and 'h'. The 'noth' feature is added below in \setmainfont. From
% https://tex.stackexchange.com/a/315635/13659 with more info here:
% http://pragma-ade.com/general/manuals/fonts-mkiv.pdf
\directlua{
  fonts.handlers.otf.addfeature{
    name = "noth",
    type = "multiple",
    data = {
      ["T_h"] = { "T", "h" },
    },
  }
}

% ~~~~~~~~~~
% Main fonts
% ~~~~~~~~~~

% Arno Pro as regular font. Use letterspaced small caps (because microtype's
% tracking is bugged).
\setmainfont{Arno Pro}[%
  Ligatures=TeX,
  Numbers=OldStyle,
  RawFeature=+noth,
  SmallCapsFeatures = {%
    LetterSpace = 5,
    Letters = SmallCaps}
]

% Scala Sans as sans font (OldStyle numerals are standard)
\setsansfont{Scala Sans Pro}[%
  Scale=MatchLowercase,
  SmallCapsFeatures = {%
    LetterSpace = 5,
    Letters = SmallCaps}
]

% Fira Mono as teletype font
\setmonofont{Fira Mono Regular}[Scale=MatchLowercase]



% ~~~~~~~~~~
% Math fonts
% ~~~~~~~~~~
% SEE COMMENTS ON MATH FONTS AT THE TOP OF THIS FILE

% Note also the nice list here: tex.stackexchange.com/questions/425098/

% Some comments on symbols
% * "2202  is (unicode) partial
% * "1D715 is (unicode) italic partial
% * "2212  is (unicode) minus
% * "002C  is (unicode) comma
% * "002F  is (unicode) slash
% * "0028  is (unicode) left paranthesis [pair with "0029]
% * "005B  is (unicode) left square bracket [pair with "005D]
% * "007B  is (unicode) left curly bracket [pair with "007D]
% * "007C  is (unicode) vertical line

% Note to self: Add colors to fonts to track usage (Colour=FF0000)

% Math fallback
\setmathfont{TeX Gyre Pagella Math}[Scale=MatchUppercase]

% Math fallback, part II
% --> Replace some symbols using another font(s) [note different scalings]

% * Integral and solar symbol is too small in TeX Gyre Pagella
\setmathfont{Cambria Math}[Scale=0.84,range={\int,\odot,"007C}]

% * Parantheses are too flimsy in TeX Gyre Pagella (and not pretty in Cambria)
\setmathfont{Asana Math}[Scale=0.90,range={"0028,"0029,"005B,"005D,"007B,"007D}]

% Use main for as much as possible
\setmathfont{Arno Pro}[range={up/{num,latin,Latin,greek,Greek},%
  \mathexclam,\mathplus,\pm,\div,\dagger,\ddagger,\minus,\mathoctothorpe,%
  \mathdollar,\mathpercent,\mathampersand,\mathquestion,\mathatsign,%
  \mathsterling,\mathyen,\euro,\increment,\infty,\checkmark,%
  \less,\equal,\greater,\ne,\leq,\geq,\matheth,\ell,"2202,"2212,"002C,"002F,},%
Script=Latin, script-features={}, sscript-features={}]

% More math families
\setmathfont{Arno Pro Italic}[range={it/{latin,Latin,greek,Greek}},%
Script=Latin, script-features={}, sscript-features={}]
\setmathfont{Arno Pro Bold}[range = {bfup/{num,latin,Latin,greek,Greek}},%
Script=Latin, script-features={}, sscript-features={}]
\setmathfont{Arno Pro Bold Italic}[range={bfit/{latin,Latin,greek,Greek}},%
Script=Latin, script-features={}, sscript-features={}]

% Set Arno Pro with the swash style as the math cal font
\newfontface\swash[Ligatures=TeX,Numbers=OldStyle,Style=Swash]{Arno Pro Italic}
\setmathfont{Arno Pro Italic}[range={cal,bfcal}, RawFeature={+swsh},
Script=Latin, script-features={}, sscript-features={}]


%%%%%%%%%%%%%%%%%%%%%%%%%%%%%%%%%%%%%%%%%%%%%%%
% Page 'n' stuff
%%%%%%%%%%%%%%%%%%%%%%%%%%%%%%%%%%%%%%%%%%%%%%%
% Page set-up
\setlrmarginsandblock{3cm}{*}{1} % Højre-venstre-marginer
\setulmarginsandblock{*}{3cm}{1} % Top-bund-marginer
\checkandfixthelayout
\sloppybottom
\strictpagechecktrue

% Basic pagestyle
\makeoddfoot{plain}{}{\sffamily\small\thepage \ / \thelastpage}{}

% Section headings
\setsecheadstyle{\large\bfseries\sffamily\raggedright}
\setsubsecheadstyle{\normalsize\bfseries\sffamily\raggedright}

% Chapter headings
\renewcommand*{\chaptitlefont}{\Large\sffamily\bfseries\raggedright}
\renewcommand*{\chapnumfont}{\chaptitlefont}    % To harmonise the font sizes
\renewcommand*{\chapnamefont}{\chaptitlefont}   % idem

% Title
\usepackage{titling}
\pretitle{\begin{center}\sffamily\LARGE}
\posttitle{\end{center}}
\preauthor{\begin{center}\sffamily\large}
\postauthor{\end{center}}
\predate{\begin{center}\sffamily}
\postdate{\end{center}}
\setlength{\droptitle}{-6em}

% Magic  --> If only one page: No header or footer
\AtEndDocument{\ifnum\value{page} = 1\thispagestyle{empty}\fi}


%%%%%%%%%%%%%%%%%%%%%%%%%%%%%%%%%%%%%%%%%%%%%%%
% Figures and floats
%%%%%%%%%%%%%%%%%%%%%%%%%%%%%%%%%%%%%%%%%%%%%%%
% Include graphics
\usepackage{graphicx}

% Width of caption
\captionwidth{.9\linewidth}
\changecaptionwidth

% Trick to automatically end captions with a period
\captiontitlefinal{.}

% Allow subfloats in figure environment
\newsubfloat{figure}

% Cool tables with footnotes
\usepackage{threeparttable, etoolbox}
\appto\TPTnoteSettings{\footnotesize}  % Requires etoolbox

% Pretty colors
\usepackage{xcolor}


%%%%%%%%%%%%%%%%%%%%%%%%%%%%%%%%%%%%%%%%%%%%%%%
% Science
%%%%%%%%%%%%%%%%%%%%%%%%%%%%%%%%%%%%%%%%%%%%%%%
% Math
\usepackage{mathtools}

% Units
\usepackage{siunitx}
\sisetup{separate-uncertainty=true}
\DeclareSIUnit\au{AU}
\DeclareSIUnit\year{yr}
\DeclareSIUnit\erg{erg}
\DeclareSIUnit\msun{M_{\odot}}
\DeclareSIUnit\lsun{L_{\odot}}

% Delimeters
\DeclarePairedDelimiter\abs{\lvert}{\rvert}
\DeclarePairedDelimiter\norm{\langle}{\rangle}


%%%%%%%%%%%%%%%%%%%%%%%%%%%%%%%%%%%%%%%%%%%%%%%
% References 'n' stuff
%%%%%%%%%%%%%%%%%%%%%%%%%%%%%%%%%%%%%%%%%%%%%%%
% Smart quotations
\usepackage[danish=guillemets]{csquotes}

% URL's
\usepackage{url}

% Nice itemizations
\usepackage{enumitem}
% \firmlists

% Smart references
\usepackage{varioref}
\usepackage{cleveref}

% Debugging
\usepackage[margin,draft]{fixme} \fxusetheme{color}

% Format of date
\usepackage[useregional]{datetime2}
\DTMlangsetup[en-GB]{ord=omit}


%%%%%%%%%%%%%%%%%%%%%%%%%%%%%%%%%%%%%%%%%%%%%%%
% Macros
%%%%%%%%%%%%%%%%%%%%%%%%%%%%%%%%%%%%%%%%%%%%%%%
% Nice spacing in macros
\usepackage{xspace}

% Small division
\newcommand\starbreak{\medskip\fancybreak{$*$\qquad$*$\qquad$*$}\medskip}

% Seperation in equation
\newcommand\eqsep{\ensuremath{\quad , \quad}}

% Lazyness
\newcommand\teff{\ensuremath{T_{\textup{eff}}}\xspace}
\newcommand\logg{\ensuremath{\log g}\xspace}

% Name
\newcommand\numpy{NumPy\xspace}
\newcommand\scipy{SciPy\xspace}
\newcommand\kepler{\textsl{Kepler}\xspace}

% Diff
\newcommand\diff[2]{\frac{\text{d}#1}{\text{d}#2}}
\newcommand\pdiff[2]{\frac{\partial #1}{\partial #2}}

% More lazy stuff
\newcommand\parname[2]{\left(#1\right)_{\text{#2}}}

% Gradients
\newcommand\radgrad{\ensuremath{\nabla_{\textup{rad}}}\xspace}
\newcommand\adgrad{\ensuremath{\nabla_{\textup{ad}}}\xspace}

% Vectors (we are doing physics!)
\renewcommand{\vec}[1]{\ensuremath{\boldsymbol{#1}}\xspace}

% Nice subscript
\newcommand\var[2]{\ensuremath{#1_{\textup{#2}}\,}\xspace}
%%%%%%%%%%%%%%
%% The end! %%
%%%%%%%%%%%%%%
